\section{Introduction}
\label{sec:intro}

\textit{Arabidopsis thaliana} grows in a light environment that switches off
every day. During the light phase the plant assimilates carbon and stores part
of it as transient leaf starch. During the night this reserve is the only carbon
source for respiration, maintenance and growth. Measurements show that the
reserve is degraded at an almost constant rate and is nearly exhausted at dawn
\cite{Graf2010,Smith2007}. The plant therefore behaves as if it rations a known
amount of carbon over a night of known length \cite{Scialdone2013}. When the
photoperiod changes, both the amount and the length change together, and the
rationing problem changes with them \cite{Sulpice2014,Stitt2012}. This makes
photoperiod a useful test case for a simulation model, because a single
environmental variable moves the light-phase carbon gain and the dark-phase
reserve use at the same time.

The problem is difficult to simulate because it couples processes on very
different time scales. Stomatal opening and leaf assimilation respond within
minutes to light, temperature and humidity. Starch and sucrose pools turn over
within hours. Nutrient status changes over days and biomass accumulates over
weeks. The carbon available for growth at any instant is not a property of any
single equation. It is the result of the environment, the gas-exchange
machinery and the internal carbohydrate and nutrient state acting together
\cite{Farquhar1980,ThornleyJohnson1990}. A model that describes only
photosynthesis, or only biomass accumulation, cannot represent this. For a
simulation study the main task is therefore to build a stable numerical
workflow in which these subsystems exchange fluxes consistently over the
day--night cycle.

Existing models resolve parts of this chain, but usually not the whole chain at
the resolution the problem needs. Crop and field-scale simulators such as
WOFOST, APSIM and DSSAT integrate growth, water and nitrogen over seasons
\cite{deWit2019,Holzworth2014,Jones2003}. For robustness and runtime they
obtain the carbon input from a light-use-efficiency term, a radiation-use
coefficient or a look-up table, and they usually advance the state on a daily
step. Soil--plant transport models and functional--structural plant models add
detail in other directions, in water transport and in architecture
respectively \cite{Janott2011,DeReffye2020,Louarn2020}. At the other end of
the scale range, leaf gas-exchange models resolve the process that the crop
models prescribe. The Farquhar--von Caemmerer--Berry biochemical model and
Jarvis-type or optimality-based stomatal formulations describe assimilation and
conductance in detail \cite{Farquhar1980,Jarvis1976,Ball1987,Medlyn2011}, and
the coupling between them has been studied extensively
\cite{Damour2010,Buckley2017}. These models are formulated and tested as
steady-state algebraic systems at a fixed carbohydrate status, so they stop at
the leaf. Few published models compute the supply--demand operating point at
every step and then carry that carbon input through a diel starch and sucrose
economy into biomass.

This paper presents such a model and reports its evaluation. The model is a
whole-plant dynamic simulator with eleven state variables, organised as six
submodels with declared inputs, outputs and parameters
(Figure~\ref{fig:overview}). S1 maps clock time to
the environmental drivers. S2 computes stomatal conductance from a
multiplicative Jarvis formulation, which acts as the supply side of gas
exchange. S3 computes net assimilation as a function of intercellular CO$_2$
from the Farquhar C3 biochemistry, which acts as the demand side. S4 closes
supply against demand. S5 holds the diel carbohydrate economy, the starch and
sucrose pools and the nitrogen and phosphorus uptake and costs. S6 converts
sucrose into leaf and root biomass. Submodels S2--S4 form a fast algebraic
subsystem, and S5--S6 form the slow differential subsystem that is integrated
in time.

The methodological content of the work sits in four places, and this is what we
expect to transfer to other coupled biophysical simulators. First, the
gas-exchange operating point is solved algebraically at runtime. Substituting
Fick's law into the branch-wise Farquhar expression gives a quadratic equation
in the intercellular CO$_2$ concentration. This quadratic is solved branch by
branch, with an explicit admissibility rule for the roots. No inner iteration
is then required inside the derivative evaluation. Second, the fast/slow decomposition
is explicit: the algebraic subsystem is eliminated exactly inside each
derivative call, so the remaining differential system is integrated with an
explicit adaptive method. Third, the implementation is verified before it is
compared with data. The checks are the residual of the algebraic closure, an
integration refinement study, and a bookkeeping check on the starch and sucrose
subsystem.
Fourth, model behaviour is analysed with a screening step followed by a
variance-based global sensitivity analysis over declared parameter ranges.

The evaluation design is chosen so that calibration and prediction stay
separate. Four coefficients are fitted and shared across all photoperiods,
where the earlier version of this model used eight coefficients fitted
separately for each photoperiod. Predictive accuracy is then measured by
holding out whole photoperiods, that is, by fitting on three light regimes and
predicting the fourth. Under this design the model reaches a held-out mean
absolute percentage error of 4.36\,\% in relative growth rate
(Figure~\ref{fig:crossval}). This value applies when one growth coefficient is
allowed to depend linearly on photoperiod. We also measure the
range of photoperiods over which the model remains usable, and we report the
boundary of that range together with the mechanism that breaks first.

The contributions of the paper are the following.
\begin{enumerate}
\item A coupled whole-plant simulation model in which Jarvis stomatal
  conductance and Farquhar C3 photosynthesis produce the carbon input at every
  step, and that input is propagated through diel starch and sucrose dynamics,
  nutrient uptake and costs, and biomass growth (Section~\ref{sec:model}).
\item An analytic runtime solution of the supply--demand closure as a quadratic
  in the intercellular CO$_2$ concentration, embedded in a fast/slow numerical
  workflow (Sections~\ref{sec:model} and~\ref{sec:implementation}).
\item A verification protocol for the implementation that is independent of any
  comparison with data, covering the algebraic closure, the time integration
  and the carbohydrate bookkeeping (Section~\ref{sec:implementation}).
\item An evaluation with four shared fitted parameters instead of a
  per-photoperiod calibration. The design uses leave-one-photoperiod-out
  cross-validation, an ablation of the carbon-input pathway, and a global
  sensitivity analysis over declared ranges
  (Sections~\ref{sec:protocol} and~\ref{sec:results}).
\item A measured validity envelope. We report the photoperiod range in which
  the diel starch balance closes, the photoperiod at which it stops closing,
  and the bound on the simulated growth response to day length
  (Sections~\ref{sec:results} and~\ref{sec:limitations}).
\end{enumerate}

The rest of the paper is organised as follows.
Section~\ref{sec:related} reviews mechanistic growth models, gas-exchange
formulations, diel carbohydrate modelling and simulation evaluation practice,
and states the gap that this work addresses. Section~\ref{sec:model} gives the
model formulation as six submodels with their states, inputs and parameters.
Section~\ref{sec:implementation} describes the numerical implementation and the
verification results. Section~\ref{sec:protocol} defines the photoperiod
scenarios, the calibration and cross-validation design, the ablation variants
and the sensitivity protocol. Section~\ref{sec:results} reports the results,
Section~\ref{sec:discussion} interprets them, Section~\ref{sec:limitations}
states what cannot be concluded, and Section~\ref{sec:conclusion} concludes.

% ============================================================
\section{Related Work}
\label{sec:related}

This section positions the model against four bodies of work. These are
mechanistic whole-plant and crop simulation, leaf gas-exchange modelling, diel
carbohydrate turnover in \textit{Arabidopsis thaliana}, and evaluation practice
in simulation studies. The aim is to locate the contribution, not to survey
each field completely.

\subsection{Mechanistic plant growth simulation models}

Mechanistic growth models link physiological processes, microclimate and
biomass change over time. The spatial and temporal scale of a model decides
where it can be applied. Spatially, models cover the leaf, the organ, the whole
plant, the canopy or the field. Temporally, they range from minutes, for
gas exchange, through hours and days, for carbon allocation, to whole seasons,
for yield. Crop simulation platforms such as APSIM, DSSAT and WOFOST are mature
tools for growth, water and nitrogen dynamics at field scale
\cite{Holzworth2014,Jones2003,deWit2019}. Recent process-based systems extend
this to detailed agronomic management and are validated against field data
\cite{BuckleyPaules2025}. Comparative studies indicate that process-based
formulations follow seasonal variation in climate better than purely empirical
ones \cite{Mamassi2024}.

Several other model families address different aspects of the same problem.
Functional--structural plant models combine architecture with physiological
modules, and they show how spatial structure and dynamic growth can be
integrated \cite{Louarn2020,Sievnen2014,DeReffye2020}. Soil--plant water
transport models describe flow from the soil through the xylem, and a
one-dimensional formulation can be as effective as a three-dimensional one at
much lower computational cost \cite{Janott2011}. Discrete approaches based on
the discrete element method are used for root--soil contact problems
\cite{Tian2024}. Optimality and energy-budget models assume that allocation
maximises a criterion such as growth or resource-use efficiency
\cite{Russo2022}. Continuous formulations describe biomass, water or vegetation
dynamics as systems of differential equations and allow formal analysis
\cite{Meron2019,ThornleyJohnson1990}. Source--sink descriptions of carbon
partitioning between shoot and root belong to the same tradition
\cite{Marcelis1996,Feller2015,Thornley1972}.

Two properties of this literature matter here. Modular organisation is common
and useful, because it lets each component be tested on its own
\cite{Sievnen2014,Zeigler2000}. The carbon input, however, is often prescribed
rather than computed. Light-use-efficiency terms, radiation-use coefficients
and look-up tables are standard in field-scale simulators, and daily time steps
are common \cite{Jones2003,deWit2019,Holzworth2014}. This separates the carbon
source from the environment and from the internal state of the plant. The model
then cannot represent the feedback between stomatal behaviour, intercellular
CO$_2$, carbohydrate status and growth.

\subsection{Photosynthesis and stomatal conductance modelling}

Photosynthesis sets the primary carbon input, so its description controls
everything downstream. For C3 plants the standard biochemical model is the
Farquhar--von Caemmerer--Berry formulation, which relates net assimilation to
intercellular CO$_2$ through Rubisco-limited carboxylation and RuBP
regeneration \cite{Farquhar1980,vonCaemmerer2000}. Its parameters, mainly
$V_{cmax}$, $J$, $\Gamma^*$, $K_c$, $K_o$ and $R_d$, and their temperature
dependence are well documented \cite{Sharkey2007,Bernacchi2001,Dubois2007}.

Assimilation also depends on the diffusive supply of CO$_2$, which stomatal
conductance controls \cite{Lawson2014}. Stomatal models therefore connect
microclimate to carbon gain. The main formulations are the empirical
multiplicative Jarvis model, the Ball--Berry family and the Medlyn optimality
model \cite{Jarvis1976,Ball1987,Medlyn2011}. The Jarvis form expresses
conductance as a product of response functions to irradiance, vapour-pressure
deficit, temperature, CO$_2$ and plant water status \cite{Jarvis1976,YU2017,Yu2021}.
It is convenient in a simulation setting because time-varying drivers from a
scenario or from measurements enter it directly. Reviews of coupled
formulations compare these alternatives and discuss where each is appropriate
\cite{Damour2010,Buckley2017}.

Coupling the two sides introduces a feedback. Intercellular CO$_2$ sets the
assimilation rate in the biochemical model, while conductance sets how much
CO$_2$ can diffuse into the leaf. The operating point is therefore defined
implicitly and must be found at every evaluation of the model
\cite{Damour2010}. Iterative solutions are widely used, and analytic solutions
of the coupled system for particular stomatal formulations are also known
\cite{Buckley2017,Kim2003}. The Farquhar and Jarvis equations used here are
inherited physiological components, not a contribution of this paper. The
contribution is the way they are solved at runtime and coupled to the
carbohydrate, nutrient and biomass states.

\subsection{Diel carbohydrate turnover and photoperiod}

Carbohydrate dynamics are the link between assimilation and growth. In the
light phase, assimilated carbon is partitioned between mobile sugars, mainly
sucrose, and transient reserves, mainly leaf starch. In the dark phase, starch
degradation supplies carbon for respiration, maintenance and growth
\cite{Stitt2012,Smith2007}. In \textit{Arabidopsis thaliana} this degradation
is closely regulated, and the reserve is almost completely used by dawn
\cite{Graf2010,Scialdone2013,Webb2015}. Mathematical models of the process
connect circadian regulation, sucrose availability and sugar signalling to
carbon partitioning and growth \cite{Feugier2013,Seki2017}. Sucrose transport
and the separation between cytosolic and vacuolar sucrose have also been
described in detail \cite{Lohaus2000,Nadwodnik2008}.

Photoperiod changes the balance between daytime accumulation and nocturnal
consumption. A short day limits the time available for assimilation and
lengthens the interval in which growth depends on stored starch. A long day
extends direct carbon input and shortens the interval of reserve use.
Experiments over contrasting short and long photoperiods show that
\textit{Arabidopsis thaliana} adjusts partitioning and growth to day length
\cite{Sulpice2014,Mengin2017,Annunziata2017}. This makes photoperiod scenarios
a suitable evaluation setting: they perturb both phases of the diel carbon
economy through one variable, and they have documented growth and carbohydrate
responses.

Growth also depends on nutrient status. Nitrogen and phosphorus affect
photosynthetic capacity, metabolism and allocation, and mechanistic models
introduce them as internal states or as limiting factors
\cite{Gastal2002,Xu2012,Barber1995,Shen2011,Carstensen2018}. Models that
resolve the diel cycle in detail, however, usually treat assimilation as given
\cite{Scialdone2013,Graf2010,Feugier2013}. The day--night carbohydrate balance
and the gas exchange that produces it are then not solved as one system.

\subsection{Verification, validation and sensitivity practice in simulation}

Evaluation of plant models normally compares simulated and reference values for
biomass, relative growth rate, and water or nutrient fluxes
\cite{Wallach2018,Holzworth2014,Jones2003}. Two distinctions from simulation
methodology are relevant here. The first is between verification and
validation: agreement with measurements does not show that the numerical
procedures inside the simulator are correct \cite{Zeigler2000,Cellier1991}.
Residual checks against analytic solutions, integration refinement studies and
conservation or bookkeeping checks address a different question from data
agreement. The second is between calibration residual and predictive error.
When a model is fitted and evaluated against the same small set of aggregate
observations, the reported error is a calibration residual. Multi-model studies
of crop simulators show how strongly reported accuracy depends on the
calibration protocol \cite{Maiorano2017,Wallach2018}.

Ablation is a further tool for attributing behaviour to model structure. It
compares variants in which a component is removed, simplified or replaced.
For a coupled model of the kind studied here, three variants are informative.
The first uses a prescribed carbon input. The second computes assimilation at
fixed conductance. The third is the full variant with environment-responsive
conductance. Such a comparison separates the effect of the coupling from the
effect of the individual equations.

Sensitivity analysis identifies which parameters control model outputs.
Variance-based methods give first-order and total-order indices and quantify
interaction structure \cite{Sobol2001,Saltelli2010}. Screening methods of the
Morris type are cheaper and indicate nonlinearity and interaction before the
full analysis \cite{Morris1991,Campolongo2007}, and Latin hypercube designs are
standard for the sampling step \cite{McKay1979}. Because these indices are
themselves Monte Carlo estimates, testing the estimator against a function with
closed-form indices is good practice \cite{Ishigami1990}. Two limits should be
stated explicitly. Indices describe model behaviour over the declared parameter
domain, and they do not establish physiological importance. Screening measures
indicate interaction, and they are not an exact variance decomposition.

\subsection{Open problems addressed here}

Read together, this literature leaves a specific gap. Leaf-level gas-exchange
models solve the supply--demand operating point accurately, but at a fixed
carbohydrate status and without downstream dynamics
\cite{Farquhar1980,Jarvis1976,vonCaemmerer2000}. Whole-plant and crop
simulators carry growth over days to seasons, but they prescribe or simplify
the carbon input and rarely resolve the diel cycle
\cite{Jones2003,deWit2019,Holzworth2014}. Dedicated diel carbohydrate models
resolve the day--night starch balance, but they take assimilation as an input
\cite{Scialdone2013,Graf2010}. The feedback loop that connects conductance,
intercellular CO$_2$, carbohydrate status and growth is therefore opened
somewhere in each family. Two further problems are methodological. The
algebraic constraint created by the coupling has to be solved at every step,
yet the solver and its behaviour outside the calibrated range are often left
unstated. And calibration is frequently performed against the same aggregate
observations used for evaluation, so a predictive error cannot be separated
from a calibration residual. Table~\ref{tab:related} places the present model
against the closest published alternatives on these axes.

This paper closes part of that gap. We do not propose new process equations for
photosynthesis, stomatal conductance or starch turnover; those are taken from
the literature. What we add is the coupled formulation with a computed carbon
input, together with its analytic runtime closure and fast/slow numerical
treatment. We also add a verification protocol that is independent of data
agreement. And we add an evaluation design in which whole photoperiods are held
out and the fitted parameter set is shared across conditions. The aim is not to replace crop simulators or
functional--structural models. It is to provide a compact and reproducible core
in which gas exchange is propagated explicitly into downstream states.

The paper does not close the rest of the gap, and we say so with numbers in
Section~\ref{sec:results}. The model does not generate the full observed
photoperiod dependence of growth from its shared dynamics; part of that
dependence remains localised in one growth-conversion coefficient. The diel
starch balance closes only up to an upper photoperiod, and above that boundary
the reserve accumulates from day to day instead of being consumed
(Figure~\ref{fig:envelope}). Canopy geometry, water transport, temperature
acclimation and multi-species parameterisation are outside the present scope.
These points are stated as named structural limitations in
Section~\ref{sec:limitations} rather than left for the reader to find.

\begin{table}[!htbp]
\centering
\caption{The present model against the closest related model families. FvCB:
Farquhar--von Caemmerer--Berry biochemistry. ``Carbon input'' distinguishes
models that compute assimilation from gas exchange at every step
(\emph{comp.}) from those that prescribe it through a light-use-efficiency
term, a radiation-use coefficient or a look-up table (\emph{presc.}).
``Fast--slow'' marks an explicit numerical separation of the algebraic
gas-exchange subsystem from the differential growth subsystem. ``Num.\ verif.''
marks published convergence, conservation or estimator-verification evidence,
as distinct from agreement with data. Entries describe the standard
formulations as given in the cited references; individual modules and versions
differ.}
\label{tab:related}
\footnotesize
\setlength{\tabcolsep}{3.5pt}
\begin{tabular}{@{}lcccccccc@{}}
\toprule
Model & FvCB & Stom. & Starch/ & N/P & Carbon & Fast-- & Num. & Open \\
      &      & model & sucrose &     & input  & slow   & verif.& src \\
\midrule
APSIM \cite{Holzworth2014}              & --  & --      & --   & N    & presc. & -- & -- & yes \\
DSSAT \cite{Jones2003}                  & --  & --      & --   & N    & presc. & -- & -- & yes \\
WOFOST \cite{deWit2019}                 & --  & --      & --   & N    & presc. & -- & -- & yes \\
GreenLab \cite{DeReffye2020}            & --  & --      & --   & --   & presc. & -- & -- & -- \\
T\&C-CROP \cite{BuckleyPaules2025}      & yes & B--B    & --   & N    & comp.  & -- & -- & yes \\
Xylem--soil \cite{Janott2011}           & --  & --      & --   & --   & presc. & -- & -- & -- \\
Scialdone et al.\ \cite{Scialdone2013}  & --  & --      & yes  & --   & presc. & -- & -- & -- \\
Graf et al.\ \cite{Graf2010}            & --  & --      & yes  & --   & presc. & -- & -- & -- \\
Feugier and Satake \cite{Feugier2013}   & --  & --      & yes  & --   & presc. & -- & -- & -- \\
Weraduwage et al.\ \cite{Weraduwage2015}& yes & --      & part.& --   & comp.  & -- & -- & -- \\
\midrule
This work                               & yes & Jarvis  & yes  & N, P & comp.  & yes & yes & yes \\
\bottomrule
\end{tabular}
\end{table}

% TODO(AUTHOR): the sentence in Section~\ref{sec:intro} that contrasts four
% shared fitted coefficients with eight coefficients fitted per photoperiod
% currently attributes the eight-parameter baseline to "the earlier version of
% this model" without a citation. Should this cite \cite{Tedone2020}, or is the
% eight-per-photoperiod calibration specific to the previous version of the
% present model and therefore left uncited?
